\documentclass[12pt,a4paper,openany,oneside]{book}

% --- Paquetes ---
\usepackage[utf8]{inputenc}
\usepackage[spanish, es-noshorthands]{babel}
\usepackage{amsmath, amsfonts, amssymb}
\usepackage{graphicx}
\usepackage{geometry}
\usepackage{hyperref}
\usepackage{booktabs}
\usepackage{fancyhdr}
\usepackage{listings}
\usepackage{xcolor}
\usepackage{tikz}
\usepackage{pgfplots}
\usepackage{colortbl}
\usepackage{multirow}
\usepackage{hhline}
\usetikzlibrary{arrows.meta, positioning, shapes.geometric, calc, shapes}
\pgfplotsset{compat=1.18}

\geometry{margin=2.5cm}

% --- Colores y estilos para código Python ---
\definecolor{codegreen}{rgb}{0,0.6,0}
\definecolor{codegray}{rgb}{0.5,0.5,0.5}
\definecolor{codepurple}{rgb}{0.58,0,0.82}
\definecolor{backcolour}{rgb}{0.95,0.95,0.92}

\lstset{
    language=Python,
    backgroundcolor=\color{backcolour},   
    commentstyle=\color{codegreen},
    keywordstyle=\color{blue},
    numberstyle=\tiny\color{codegray},
    stringstyle=\color{codepurple},
    basicstyle=\ttfamily\small,
    breaklines=true,
    frame=single,
    showstringspaces=false
}

% --- Estilo de cabecera y pie ---
\pagestyle{fancy}
\fancyhf{}
\renewcommand{\headrulewidth}{0.4pt}
\renewcommand{\footrulewidth}{0.4pt}
\fancyhead[L]{\small Carlos César Sánchez Coronel}
\fancyhead[R]{\small Sesión 1: Introducción a NLP}
\fancyfoot[L]{\small Carlos César Sánchez Coronel}
\fancyfoot[C]{\thepage}

% --- Portada ---
\title{
    \textbf{Sesión 1} \\ 
    \Huge \textbf{INTRODUCCIÓN A NLP Y FUNDAMENTOS LINGÜÍSTICOS} \\
    \Large Del lenguaje humano a su procesamiento computacional
}
\author{
    \small Por \\ \vspace{0.2cm}
    \Large \textsc{Carlos César Sánchez Coronel}
}
\date{2026}

\begin{document}

\maketitle
\tableofcontents
\addcontentsline{toc}{chapter}{Índice general}

% ------------------------------------------------------------------
% CAPÍTULO 1: SESIÓN 1 - INTRODUCCIÓN A NLP
% ------------------------------------------------------------------
\chapter{Introducción a NLP y Fundamentos Lingüísticos}

En esta primera sesión se establecen las bases conceptuales del Procesamiento del Lenguaje Natural (NLP). Se exploran los niveles del lenguaje, las tareas fundamentales como el etiquetado gramatical y el análisis sintáctico, y se examinan corpus anotados de diversos dominios. El objetivo es comprender la complejidad del lenguaje humano y cómo la lingüística computacional aborda estos desafíos.

\section{Logro de la sesión}
Comprender el alcance del NLP, sus aplicaciones en la industria y los conceptos lingüísticos básicos que subyacen al procesamiento automático del lenguaje humano.

\section{Introducción: ¿Qué es NLP y por qué es importante?}

\subsection{Definición y objetivos}
El \textbf{Procesamiento del Lenguaje Natural (NLP)} es una subdisciplina de la inteligencia artificial que se ocupa de dotar a las computadoras de la capacidad de entender, interpretar y generar lenguaje humano de manera útil y significativa.

A diferencia de otros campos de la IA que trabajan con datos numéricos o categóricos, el NLP enfrenta desafíos únicos:
\begin{itemize}
    \item \textbf{Ambigüedad}: una misma palabra puede tener múltiples significados según el contexto.
    \item \textbf{Creatividad}: el lenguaje evoluciona constantemente, aparecen nuevas palabras y expresiones.
    \item \textbf{Dependencia del contexto}: el significado de una oración puede depender de oraciones anteriores.
    \item \textbf{Conocimiento de mundo}: muchas expresiones requieren conocimiento extralingüístico.
\end{itemize}

\subsection{Breve historia del NLP}
\begin{itemize}
    \item \textbf{1950s-1960s}: enfoques basados en reglas y sistemas expertos (traducción automática palabra por palabra).
    \item \textbf{1980s-1990s}: auge de los métodos estadísticos y corpus anotados.
    \item \textbf{2000s}: modelos de machine learning clásicos aplicados a texto (Naive Bayes, SVM).
    \item \textbf{2010s}: word embeddings (Word2Vec, GloVe) y redes neuronales.
    \item \textbf{2017-presente}: Transformers, BERT, GPT, modelos fundacionales.
\end{itemize}

\subsection{El desafío fundamental: del lenguaje humano a representación computacional}
El lenguaje humano es inherentemente discreto, simbólico y estructurado, mientras que las computadoras operan con números. El NLP consiste en encontrar transformaciones que preserven el significado y la estructura mientras permiten el cómputo eficiente.

\section{Niveles del lenguaje: de lo más simple a lo más complejo}

\subsection{Morfología: la estructura de las palabras}
\textbf{Definición formal}: La morfología estudia la estructura interna de las palabras y sus constituyentes mínimos con significado, llamados \textbf{morfemas}.

\textbf{Ejemplo práctico}: La palabra "inextinguible" se descompone en:
\begin{itemize}
    \item Prefijo "in-" (negación)
    \item Lexema "extingu" (raíz que aporta el significado básico)
    \item Sufijo "-ible" (posibilidad)
\end{itemize}

\textbf{Problemas computacionales}:
\begin{itemize}
    \item \textbf{Segmentación morfológica}: dividir palabras en morfemas.
    \item \textbf{Lematización}: reducir una palabra a su forma canónica (ej. "corrió", "correrán" → "correr").
    \item \textbf{Stemming}: reducir a una raíz aproximada (algoritmos como Porter Stemmer).
\end{itemize}

\textbf{Aplicaciones}:
\begin{itemize}
    \item Búsqueda de información: no importa si el usuario busca "corriendo", "corredor" o "correr".
    \item Análisis de sentimiento: prefijos negativos cambian la polaridad ("feliz" vs "infeliz").
\end{itemize}

\subsection{Sintaxis: la estructura de las oraciones}
\textbf{Definición formal}: La sintaxis estudia cómo se combinan las palabras para formar oraciones gramaticales, definiendo reglas de orden y concordancia.

\textbf{Estructura de constituyentes}: Una oración puede descomponerse jerárquicamente:
\begin{itemize}
    \item Oración → Sintagma Nominal + Sintagma Verbal
    \item Sintagma Nominal → (Determinante) + Nombre + (Adjetivo)
    \item Sintagma Verbal → Verbo + Sintagma Nominal
\end{itemize}

\textbf{Representación matemática}: Los árboles sintácticos son la base del \textbf{parsing} o análisis sintáctico. Formalmente, un árbol sintáctico es un grafo dirigido acíclico donde:
\begin{itemize}
    \item Los nodos representan constituyentes (SN, SV, N, V, etc.)
    \item Las hojas son las palabras
    \item Las relaciones padre-hijo representan la composición
\end{itemize}

\textbf{Ejemplo}: La oración "El gato persigue al ratón" tiene un árbol donde "El gato" es un SN sujeto, "persigue al ratón" es un SV predicado, y "al ratón" es un SN objeto.

\subsection{Semántica: el significado}
\textbf{Definición formal}: La semántica estudia el significado de las palabras, las oraciones y las relaciones entre ellas.

\textbf{Niveles semánticos}:
\begin{itemize}
    \item \textbf{Semántica léxica}: significado de palabras individuales y relaciones entre ellas (sinonimia, antonimia, hiperonimia).
    \item \textbf{Semántica composicional}: cómo se combina el significado de las palabras para formar el significado de la oración.
    \item \textbf{Semántica del discurso}: significado en contexto extendido.
\end{itemize}

\textbf{Problemas computacionales}:
\begin{itemize}
    \item \textbf{Desambiguación del sentido de las palabras (WSD)}: determinar qué acepción de una palabra polisémica se está utilizando. Por ejemplo, "banco" puede ser institución financiera, asiento o conjunto de peces.
    \item \textbf{Reconocimiento de entidades nombradas (NER)}: identificar nombres de personas, organizaciones, lugares, fechas, etc.
    \item \textbf{Análisis de relaciones semánticas}: detectar relaciones como "trabaja\_en" entre una persona y una organización.
\end{itemize}

\textbf{Corpus anotado semánticamente}: Un ejemplo relevante es el proyecto \textbf{ELEXIS}, que ha producido corpus de 2 millones de palabras con anotación de desambiguación semántica para los idiomas oficiales de la Unión Europea. Estos corpus incluyen identificadores de synsets de BabelNet y WordNet, permitiendo entrenar modelos que determinan el significado correcto de las palabras en contexto.

\subsection{Pragmática: el uso en contexto}
\textbf{Definición formal}: La pragmática estudia cómo el contexto influye en la interpretación del significado, incluyendo la intención del hablante, el conocimiento compartido y las normas sociales.

\textbf{Ejemplos}:
\begin{itemize}
    \item Ironía: "Qué buen tiempo hace" durante una tormenta.
    \item Implicaturas: "¿Tienes hora?" no es una pregunta sobre posesión, sino una solicitud de información.
    \item Actos de habla: prometer, preguntar, ordenar, etc.
\end{itemize}

\textbf{Desafíos computacionales}:
\begin{itemize}
    \item Detección de sarcasmo/ironía en redes sociales.
    \item Comprensión de intenciones en asistentes virtuales.
\end{itemize}

\subsection{Interrelación de los niveles}
Ningún nivel opera de forma aislada:
\begin{itemize}
    \item La sintaxis necesita la morfología (concordancia de género y número).
    \item La semántica necesita la sintaxis (el significado depende de quién hizo qué a quién).
    \item La pragmática necesita la semántica (la ironía requiere entender el significado literal primero).
\end{itemize}

\section{Part-of-Speech (POS) Tagging}

\subsection{Definición}
El etiquetado de partes de la oración consiste en asignar a cada palabra de un texto su categoría gramatical (sustantivo, verbo, adjetivo, etc.) de acuerdo con su contexto.

\subsection{Conjuntos de etiquetas (tagsets)}
Existen diferentes estándares, siendo los más comunes:
\begin{itemize}
    \item \textbf{Penn Treebank tagset} (inglés): 45 etiquetas como NN (sustantivo singular), NNS (sustantivo plural), VB (verbo base), VBD (verbo pasado), JJ (adjetivo), RB (adverbio).
    \item \textbf{EAGLES} (para lenguas europeas): estándar más detallado.
\end{itemize}

\subsection{Formalización matemática del problema}
Dada una secuencia de palabras \( w_1, w_2, ..., w_n \), queremos encontrar la secuencia de etiquetas \( t_1, t_2, ..., t_n \) más probable:

\[
\hat{t}_1^n = \arg\max_{t_1^n} P(t_1^n | w_1^n)
\]

Por el teorema de Bayes:

\[
P(t_1^n | w_1^n) = \frac{P(w_1^n | t_1^n) P(t_1^n)}{P(w_1^n)}
\]

Como \( P(w_1^n) \) es constante para todas las secuencias de etiquetas, maximizar lo anterior equivale a maximizar \( P(w_1^n | t_1^n) P(t_1^n) \).

\subsection{Modelos de etiquetado}
\begin{itemize}
    \item \textbf{Modelos basados en reglas}: usando contexto lingüístico (ej. "si palabra termina en -ción y es singular, etiquetar como NN").
    \item \textbf{Modelos estadísticos}: como el \textbf{Hidden Markov Model (HMM)}, que asume:
    \[
    P(t_1^n) \approx \prod_{i=1}^n P(t_i | t_{i-1}), \quad
    P(w_1^n | t_1^n) \approx \prod_{i=1}^n P(w_i | t_i)
    \]
    El entrenamiento estima las probabilidades de transición \( P(t_i | t_{i-1}) \) y de emisión \( P(w_i | t_i) \) a partir de corpus anotados.
    \item \textbf{Modelos basados en aprendizaje profundo}: redes neuronales (BiLSTM, Transformers) que capturan contexto a ambos lados de la palabra.
\end{itemize}

\subsection{Aplicaciones}
El POS tagging es un paso fundamental para:
\begin{itemize}
    \item Parsing sintáctico.
    \item Reconocimiento de entidades nombradas.
    \item Análisis de sentimiento (los adjetivos suelen cargar la polaridad).
\end{itemize}

\subsection{Corpus anotado con POS}
El proyecto \textbf{Petro NLP} para la industria del petróleo y gas ha desarrollado el corpus \textbf{Petrolês}, que incluye anotación de etiquetas morfosintácticas (POS) en documentos técnicos geocientíficos en portugués. Este corpus fue creado por un equipo interdisciplinario de geocientíficos, lingüistas e ingenieros, demostrando cómo los recursos lingüísticos deben adaptarse al dominio específico.

\section{Parsing: Análisis Sintáctico}

\subsection{Definición}
El parsing es el proceso de analizar una oración para determinar su estructura sintáctica de acuerdo con una gramática formal.

\subsection{Tipos de parsing}

\subsubsection{Parsing de constituyentes (phrase structure parsing)}
Construye un árbol que muestra cómo las palabras se agrupan en frases y cómo estas frases se combinan para formar la oración.

\textbf{Ejemplo}: Para "El gato persigue al ratón", un árbol de constituyentes tendría:
\begin{itemize}
    \item S (oración)
    \begin{itemize}
        \item NP (sintagma nominal): "El gato"
        \item VP (sintagma verbal): "persigue al ratón"
        \begin{itemize}
            \item V (verbo): "persigue"
            \item NP: "al ratón"
        \end{itemize}
    \end{itemize}
\end{itemize}

\textbf{Formalismo gramatical}: Gramáticas libres de contexto (CFG) con reglas como:
\begin{align*}
    S &\to NP \; VP \\
    NP &\to Det \; N \\
    VP &\to V \; NP
\end{align*}

\textbf{Parsing probabilístico (PCFG)}: Asigna probabilidades a las reglas para resolver ambigüedades.

\subsubsection{Parsing de dependencias (dependency parsing)}
Construye un grafo donde las palabras se conectan mediante relaciones binarias (dependencias). Cada palabra tiene un "cabeza" (gobernador) y una relación.

\textbf{Ejemplo}: En la misma oración:
\begin{itemize}
    \item "persigue" es la raíz.
    \item "gato" depende de "persigue" con relación "nsubj" (sujeto nominal).
    \item "ratón" depende de "persigue" con relación "obj" (objeto).
    \item "El" depende de "gato" con relación "det" (determinante).
    \item "al" (contracción de "a el") podría analizarse como "ratón" → "case".
\end{itemize}

\textbf{Ventajas}:
\begin{itemize}
    \item Más cercano a la semántica (las relaciones son explícitas).
    \item Funciona bien para lenguas con orden libre.
    \item Más fácil de anotar (solo relaciones binarias).
\end{itemize}

\subsection{Formalización matemática del parsing de dependencias}
Un árbol de dependencias para una oración de \(n\) palabras es un grafo dirigido \( G = (V, E) \) donde:
\begin{itemize}
    \item \( V = \{1, 2, ..., n\} \) son las posiciones de las palabras.
    \item Cada vértice tiene exactamente un arco entrante, excepto la raíz (que tiene cero).
    \item El grafo es acíclico.
\end{itemize}

El objetivo es encontrar el árbol \(T\) que maximice la puntuación:
\[
T^* = \arg\max_{T} \sum_{(i,j) \in T} s(i,j)
\]
donde \(s(i,j)\) es la puntuación de que la palabra \(j\) dependa de la palabra \(i\) (con alguna relación).

\subsection{Corpus de árboles sintácticos (Treebanks)}
Un \textbf{treebank} es un corpus anotado con estructuras sintácticas, esencial para entrenar parsers estadísticos y neuronales.

\textbf{Ejemplo destacado}: El \textbf{CINTIL-TreeBank} es un corpus de árboles de constituencia para el portugués, compuesto por 10,039 oraciones (110,166 tokens) provenientes de noticias, novelas y textos de prueba. Fue creado mediante anotación semi-automática con doble ciego y adjudicación, garantizando alta calidad. Su objetivo fue proporcionar un conjunto de datos de alta calidad para desarrollar recursos y herramientas automáticas de NLP para el portugués.

\section{Corpus y Datasets Anotados}

\subsection{Definición y tipos de corpus}
Un \textbf{corpus} (plural: corpora) es una colección estructurada de textos, generalmente almacenada electrónicamente, que sirve como base para la investigación lingüística y el entrenamiento de modelos de NLP.

\textbf{Tipos según su propósito}:
\begin{itemize}
    \item \textbf{Corpus generales}: representan el lenguaje en su conjunto (ej. British National Corpus).
    \item \textbf{Corpus de dominio específico}: textos de un área particular (medicina, derecho, petróleo).
    \item \textbf{Corpus monolingües vs. multilingües}.
    \item \textbf{Corpus anotados vs. no anotados}: los anotados contienen información lingüística añadida (POS, árboles, entidades, etc.).
\end{itemize}

\subsection{Ejemplos de corpus anotados en la investigación actual}

\subsubsection{Corpus para reconocimiento de entidades médicas: RareDis}
El corpus \textbf{RareDis} fue creado para reconocer enfermedades raras y sus manifestaciones clínicas en textos médicos. Contiene 1,041 textos de Orphanet, con más de 5,000 enfermedades raras y 6,000 manifestaciones clínicas anotadas. La calidad del corpus se midió mediante acuerdo entre anotadores, alcanzando una F1 de 83.5\%. El dataset está particionado en entrenamiento, validación y prueba, listo para entrenar modelos de deep learning.

\subsubsection{Corpus de diálogos para resumen abstractivo: SAMSum}
El corpus \textbf{SAMSum} contiene más de 16,000 diálogos de chat con resúmenes escritos por humanos. Cada ejemplo tiene dos campos: `dialogue` (el texto del diálogo) y `summary` (el resumen). Es un recurso fundamental para entrenar modelos de resumen abstractivo de conversaciones.

\subsubsection{Corpus multilingüe masivo: ELEXIS}
Los \textbf{ELEXIS corpora} son una colección de 24 corpus correspondientes a los idiomas oficiales de la UE, cada uno con un tamaño objetivo de 1,000 millones de palabras. Se construyeron utilizando tecnología especializada en recopilar contenido web lingüísticamente valioso. Además, incluyen muestras de 2 millones de palabras con anotación semántica (desambiguación de sentidos).

\subsubsection{Corpus para análisis narrativo: Text2Story Lusa}
El corpus \textbf{Text2Story Lusa} se centra en artículos periodísticos en portugués europeo con anotación de componentes narrativos. Incluye 357 artículos, con un subconjunto de 117 artículos densamente anotados (más de 50,000 anotaciones individuales). Facilita la investigación en extracción de narrativas y estructuras discursivas.

\subsubsection{Corpus de dominio industrial: Petro NLP}
El proyecto \textbf{Petro NLP} ha desarrollado múltiples recursos para la industria del petróleo y gas en portugués:
\begin{itemize}
    \item \textbf{Petro KGraph}: un grafo de conocimiento poblado con entidades y relaciones del dominio.
    \item \textbf{Petrolês}: corpus de texto crudo.
    \item \textbf{PetroGold}: treebank de dependencias sintácticas.
    \item \textbf{PetroNER}: corpus anotado con entidades nombradas.
    \item \textbf{PetroRE}: corpus anotado con relaciones semánticas.
\end{itemize}

\subsubsection{Corpus para comprensión del lenguaje en catalán: NLUCat}
El dataset \textbf{NLUCat} contiene cerca de 12,000 instrucciones en catalán anotadas con intenciones (intents) y segmentos (spans). Está diseñado para asistentes virtuales domésticos, incluyendo necesidades de personas vulnerables (recordatorios de medicación, trámites administrativos). Cada ejemplo incluye la instrucción, la intención y los segmentos etiquetados con su tipo de información.

\subsection{Principios FAIR en la creación de corpus}
Los corpus modernos se adhieren a los principios \textbf{FAIR}:
\begin{itemize}
    \item \textbf{Findable}: con metadatos ricos y identificadores persistentes.
    \item \textbf{Accessible}: acceso abierto o controlado mediante protocolos estándar.
    \item \textbf{Interoperable}: formatos comunes (JSON, TSV) y vocabularios compartidos.
    \item \textbf{Reusable}: licencias claras (ej. CC BY 4.0) y documentación detallada.
\end{itemize}

\section{Aplicaciones del NLP en la industria}

\subsection{Chatbots y asistentes virtuales}
\begin{itemize}
    \item \textbf{Ejemplo}: asistentes bancarios que responden preguntas sobre saldos, transferencias.
    \item \textbf{Componentes}: clasificación de intenciones, extracción de entidades (fechas, montos), gestión de diálogo.
    \item \textbf{Caso NLUCat}: desarrollado específicamente para un asistente virtual en catalán, considerando la realidad cultural de los hablantes.
\end{itemize}

\subsection{Análisis de sentimiento y monitoreo de redes sociales}
\begin{itemize}
    \item \textbf{Ejemplo}: marcas que monitorean Twitter para detectar crisis de reputación.
    \item \textbf{Técnicas}: clasificación de polaridad (positivo, negativo, neutro), detección de emociones.
    \item \textbf{Desafío}: ironía y sarcasmo (nivel pragmático).
\end{itemize}

\subsection{Traducción automática}
\begin{itemize}
    \item \textbf{Ejemplo}: Google Translate, DeepL.
    \item \textbf{Modelos}: secuencia a secuencia con atención, Transformers.
    \item \textbf{Desafío}: preservar significado y fluidez.
\end{itemize}

\subsection{Búsqueda semántica y motores de respuesta}
\begin{itemize}
    \item \textbf{Ejemplo}: buscadores que entienden la intención detrás de las palabras.
    \item \textbf{Técnicas}: embeddings de oraciones, recuperación densa.
\end{itemize}

\subsection{Extracción de información en dominios específicos}
\begin{itemize}
    \item \textbf{Ejemplo en petróleo}: Petro NLP extrae automáticamente entidades geocientíficas de informes técnicos antiguos. La industria del petróleo tiene el problema de que el 80\% de los nuevos datos son no estructurados, y solo en el Reino Unido, el 76\% de los datos liberados eran archivos no estructurados (TIFF, PDF, Word, Excel).
    \item \textbf{Ejemplo en salud}: RareDis extrae enfermedades raras y síntomas de literatura biomédica para ayudar en diagnóstico.
\end{itemize}

\subsection{Resumen automático}
\begin{itemize}
    \item \textbf{Ejemplo}: SAMSum para resumir conversaciones de chat.
    \item \textbf{Tipos}: extractivo (selecciona frases importantes) vs. abstractivo (genera texto nuevo).
\end{itemize}

\subsection{Análisis de contratos y documentos legales}
\begin{itemize}
    \item \textbf{Ejemplo}: detección de cláusulas abusivas, extracción de fechas y obligaciones.
\end{itemize}

\section{Desafíos actuales en NLP}

\subsection{Ambigüedad y polisemia}
\begin{itemize}
    \item \textbf{Problema}: palabras con múltiples significados.
    \item \textbf{Solución}: modelos contextuales como BERT que usan el contexto para representar la palabra.
\end{itemize}

\subsection{Dependencia del contexto}
\begin{itemize}
    \item \textbf{Problema}: oraciones como "El banco me negó el préstamo" y "Me senté en el banco del parque".
    \item \textbf{Solución}: word embeddings contextuales.
\end{itemize}

\subsection{Conocimiento de mundo}
\begin{itemize}
    \item \textbf{Problema}: "El Papa visitó Barcelona. Es la primera vez que visita la ciudad." Requiere saber que el Papa es una persona, que Barcelona es una ciudad.
    \item \textbf{Solución}: integración con bases de conocimiento (grafos de conocimiento).
\end{itemize}

\subsection{Escasez de recursos para lenguas minoritarias}
\begin{itemize}
    \item \textbf{Problema}: la mayoría de los modelos se entrenan en inglés.
    \item \textbf{Solución}: transfer learning, creación de corpus como NLUCat para catalán.
\end{itemize}

\subsection{Sesgos y equidad}
\begin{itemize}
    \item \textbf{Problema}: los modelos aprenden sesgos presentes en los datos de entrenamiento.
    \item \textbf{Solución}: curado cuidadoso de datos, métricas de equidad.
\end{itemize}

\section{Preparación para entrevistas: preguntas típicas}

\subsection{Preguntas conceptuales}
\begin{itemize}
    \item \textbf{"¿Qué es NLP y cuáles son sus principales desafíos?"}
    Esperan mencionar la ambigüedad, la complejidad del lenguaje, la necesidad de conocimiento de mundo.
    \item \textbf{"Explica la diferencia entre sintaxis y semántica con ejemplos."}
    Sintaxis: estructura gramatical ("el gato persigue al ratón" vs "el ratón persigue al gato" tienen la misma estructura pero diferente significado). Semántica: significado de las palabras y su combinación.
    \item \textbf{"¿Qué es POS tagging y por qué es importante?"}
    Definición formal y aplicaciones (parsing, NER, análisis de sentimiento).
    \item \textbf{"¿Qué es un corpus anotado? Menciona ejemplos."}
    Colección de textos con anotaciones lingüísticas. Ejemplos: RareDis (enfermedades raras), Petro NLP (petróleo), SAMSum (diálogos resumidos).
\end{itemize}

\subsection{Preguntas sobre aplicaciones}
\begin{itemize}
    \item \textbf{"¿Cómo diseñarías un chatbot para atención al cliente?"}
    Clasificación de intenciones, extracción de entidades, gestión de diálogo, integración con base de conocimiento.
    \item \textbf{"¿Cómo detectarías automáticamente menciones de enfermedades en textos médicos?"}
    Usar un modelo de NER entrenado en un corpus como RareDis.
\end{itemize}

\subsection{Preguntas sobre desambiguación}
\begin{itemize}
    \item \textbf{"¿Cómo resolvería la ambigüedad de la palabra 'banco' en un sistema de NLP?"}
    Usar contexto (palabras vecinas) con modelos contextuales como BERT, o usar desambiguación basada en WordNet si el dominio es conocido.
\end{itemize}

\subsection{Preguntas sobre recursos lingüísticos}
\begin{itemize}
    \item \textbf{"¿Qué consideraciones hay al crear un corpus para un dominio específico?"}
    Involucrar expertos del dominio, definir guías de anotación, medir acuerdo entre anotadores, seguir principios FAIR.
\end{itemize}



\section{Resumen de la sesión}
\begin{itemize}
    \item El \textbf{NLP} es un campo multidisciplinario que combina lingüística, computación y matemáticas para procesar el lenguaje humano.
    \item Los \textbf{niveles del lenguaje} (morfología, sintaxis, semántica, pragmática) proporcionan un marco para entender la complejidad del problema.
    \item \textbf{POS tagging} asigna categorías gramaticales a las palabras, fundamental para análisis posteriores.
    \item \textbf{Parsing} descubre la estructura sintáctica (constituyentes o dependencias), esencial para entender relaciones.
    \item Los \textbf{corpus anotados} son la base del NLP moderno; ejemplos como RareDis, SAMSum, Petro NLP, ELEXIS y NLUCat ilustran la diversidad de dominios y tipos de anotación.
    \item Las \textbf{aplicaciones industriales} abarcan desde chatbots hasta extracción de información en dominios específicos como petróleo y salud.
    \item Los \textbf{desafíos actuales} incluyen ambigüedad, dependencia del contexto, escasez de recursos y sesgos.
\end{itemize}

\end{document}